Tato příručka shrnuje principy a konkrétní nástroje pro zodpovědné a efektivní začlenění generativní AI do výuky na VUT. Jejím cílem není poskytovat vyčerpávající technický manuál, ale nabídnout praktický rámec kroků, které lze okamžitě aplikovat v kurzech i na úrovni kateder. Zaměřuje se na principy bezpečného nasazení, ověřitelnosti výsledků a minimalizace rizik při zachování didaktické hodnoty výuky. Hlavní sdělení, která si ponechte jako praktický rámec pro další kroky, jsou tato:

- AI je prostředek pro zlepšení kvality výuky — především pro personalizaci, rychlejší zpětnou vazbu a tvorbu variabilních cvičení — nikoli náhrada didaktické práce učitele. Role pedagoga zůstává v navrhování výukových cílů, ověřování výsledků učení a rozhodování o etickém rámci použití nástrojů (general background knowledge). Pedagogové také určují, které části úlohy mohou být delegovány na nástroj a které vždy musí projít lidskou kontrolou.
- Začněte malými, měřitelnými piloty a ověřujte dopady na učení i integritu hodnocení. Krokové piloty minimalizují rizika a umožňují iterativní zlepšování na základě konkrétních metrik a empirických pozorování (general background knowledge). Pro piloty doporučujeme stanovit kontrolní skupinu nebo předpilotní baseline, aby byly výsledky porovnatelné a interpretovatelné.

Konkrétní doporučení a ukázky, které najdete v příručce, jsou zaměřené na rychlé nasazení v hodině i na úrovni katedry. Například použití vzdělávacího prostředí Minecraft Education jako demonstračního nástroje pro vysvětlení principů strojového učení a základních konceptů AI je ověřený praktický přístup, který lze rychle implementovat ve výuce \cite{kb_ai_101_tipu_pdf_04e4a115_page_15_2}. Podobně nástroje z rodiny Microsoft (funkce „Pokrok v matematice“ či „Pokrok ve čtení“) ilustrují, jak lze AI využít pro generování cvičení, sběr postupů a sledování individuálního pokroku žáků/studentů \cite{kb_ai_101_tipu_pdf_04e4a115_page_8_1, kb_ai_101_tipu_pdf_04e4a115_page_7_1}. V příručce jsou k těmto případům doplněny konkrétní lekce, prompt‑šablony a návrhy metrik ke sběru dat.

Další kroky — doporučené krátké akce (praktický checklist)
\begin{itemize}
\item Vyberte malý pilotní kurz (1–2 týmy/sekce) a definujte jasné metriky úspěchu (zapojení studentů, kvalita zpětné vazby, čas vyučujících). Stanovte také časový rámec pilotu a zodpovědné osoby za sběr dat a vyhodnocení. (general background knowledge) Doporučujeme vymezit minimální a cílové hodnoty metrik, aby bylo možné rozhodnout o rozšíření projektu.
\item Použijte šablony a prompt‑knihovnu z repozitáře pro rychlé vytvoření testovacích aktivit a rubrik; validujte každý generovaný materiál lidskou kontrolou před nasazením. Doporučujeme záznam kontroly (kdo, kdy, jaké změny) pro zajištění reprodukovatelnosti a odpovědnosti. (general background knowledge) Přidejte do repozitáře verze šablon a log změn pro auditní stopu.
\item Zajistěte transparentnost: upravte sylaby tak, aby obsahovaly pravidla pro používání AI a povinnost deklarace (viz kap. 2 a přílohy). Do sylabu uveďte také konkrétní příklady povoleného a zakázaného použití, aby studenti přesně věděli, jak jednat. (general background knowledge) Doporučujeme rovněž formulovat sankce a kroky v případě porušení těchto pravidel.
\item Pro piloty využijte ilustrativní nástroje (např. Minecraft Education pro demonstrace, Microsoft funkce pro procvičování) podle doporučení v kapitolách a videonávodech \cite{kb_ai_101_tipu_pdf_04e4a115_page_15_2, kb_ai_101_tipu_pdf_04e4a115_page_8_1}. Využijte ukázkové lekce a připravené scénáře z repozitáře pro rychlejší start a srovnání výsledků. Přizpůsobte scénáře konkrétním učebním cílům a věkové skupinám.
\item Dokumentujte DPIA/DPA a bezpečnostní opatření tam, kde jsou zpracovávána osobní data studentů; zapojte IT/právní oddělení k revizi. Ujistěte se o minimalizaci sdílených osobních údajů a o šifrování či jiných technických opatřeních tam, kde to je nutné. (general background knowledge) V případě potřeby zařiďte informovaný souhlas studentů nebo jejich zákonných zástupců.
\item Monitorujte etické aspekty a bias v datech: zaznamenávejte známé omezení nástrojů a informujte o nich studenty i kolegy, aby bylo používání transparentní.
\end{itemize}

Praktické provozní doporučení pro piloty
\begin{itemize}
\item Zaveďte krátké školení pro vyučující a asistenty o fungování používaných nástrojů, správném promptování a základních rizicích (bias, hallucinations). Součástí školení by měla být i praktická část s ukázkami ověřování výstupů.
\item Určete jasné postupy pro ověřování autenticity studentské práce a pro řešení podezření z porušení akademických pravidel při použití AI. Doporučujeme kombinovat techniky (plagiátorský software, ústní obhajoby, laboratorní záznamy) pro robustní ověření.
\item Sledujte časovou náročnost: měřte, zda AI řešení skutečně šetří čas vyučujících, a zvažte celkovou režii související s validací výstupů. Zaznamenávejte skutečné časy na přípravu, korekci a správu aktivit.
\item Zaznamenávejte zpětnou vazbu studentů systematicky, aby bylo možné upravovat prompt šablony a aktivity na základě uživatelských zkušeností. Využijte krátké anonymní ankety i fokusové rozhovory pro hlubší vhled.
\item Zaveďte malé kontrolní mechanismy kvality: náhodné audity generovaných materiálů, peer review mezi vyučujícími a běžné aktualizace prompt‑knihovny.
\end{itemize}

Pozvánka k participaci a udržování praxe
(general background knowledge) Tento materiál je navržen jako „living document“ — sdílejte svoji zkušenost, přispívejte šablonami a případovými studiemi, a účastněte se kvartálních aktualizací repozitáře a video‑tutoriálů. Zapojení širší komunity vyučujících urychlí vznik ověřených vzorů a pomůže rozšířit osvědčené postupy napříč fakultami. Navrhujeme také pořádat interní workshopy pro předávání znalostí mezi katedrami a pravidelné evaluační setkání po ukončení každého pilotu. Do repozitáře ukládejte nejen úspěšné případy, ale i neúspěchy a klíčová rizika, aby ostatní mohli těžit z vaší zkušenosti.

Závěrem: experimentujte zodpovědně, měřte výsledky a sdílejte poznatky. AI může posunout kvalitu výuky, pokud zůstane nástrojem pedagoga — nikoli náhradou. Dodržováním zásad transparentnosti, ochrany dat a pečlivého ověřování výstupů lze minimalizovat rizika a maximalizovat přínos pro studenty i vyučující. Udržujte dialog v rámci komunity, aktualizujte praktiky podle nových poznatků a považujte piloty za iterativní proces učení pro všechny zúčastněné.